% Full presentation deck: one slide block per person, in the order of the report.
% P2 assembles it. Compile from this folder: pdflatex slides.tex
\documentclass[aspectratio=169,11pt]{beamer}
\usetheme{default}
\setbeamertemplate{navigation symbols}{}
\setbeamertemplate{footline}[frame number]
\usepackage{amsmath,amssymb,booktabs}

\title{Optimal Stopping with Memory}
\subtitle{Deep recurrent policies for American options under rough volatility\\
and multi-market backtesting}
\author{Anamay Mittal, Vraj Shah, Anshul D Jain, Sachit Kapoor, Mridul Kwatra}
\date{MAC F244 Stochastic Calculus}

\begin{document}

\begin{frame}
\titlepage
\end{frame}

% Slide block for P2 (theory and baselines) and P1 (OSPG and replication).
% Frames only; included by slides.tex.

% =============================================================== P2: theory
\section{Optimal stopping}

\begin{frame}{American options are optimal stopping problems}
\begin{itemize}
\item Exercising at date $j$ pays $g_j=e^{-rt_j}(K-S_j)^+$; the price is
  $V_0=\sup_{\tau}\mathbb E^{\mathbb Q}[g_\tau]$ over stopping times $\tau$
  ($\{\tau=j\}\in\mathcal F_j$: decide using the past only).
\item \textbf{Snell envelope:} $U_H=g_H$, \ $U_j=\max\big(g_j,\ \mathbb E[U_{j+1}\mid\mathcal F_j]\big)$.
  \begin{itemize}
  \item smallest supermartingale above $g$; \ $\tau^*=\min\{j:U_j=g_j\}$, \ $V_0=U_0$
  \end{itemize}
\item Exercise when the payoff beats the \emph{continuation value} $C_j=\mathbb E[U_{j+1}\mid\mathcal F_j]$.
\item Markov prices: $C_j=c_j(S_j)$, so tree, Longstaff--Schwartz and Deep Optimal Stopping
  only need today's state. \textbf{With memory (rough volatility) they do not.}
\end{itemize}
\end{frame}

\begin{frame}{Continuous time: a free-boundary problem}
\begin{columns}[T]
\begin{column}{0.55\textwidth}
\begin{itemize}
\item $\min\big(-\partial_tV-\mathcal LV,\ V-g\big)=0$
  \begin{itemize}
  \item continue: Black--Scholes / Feynman--Kac PDE
  \item exercise: $V=K-S$
  \item boundary $L(t)$: smooth pasting $\partial_SV=-1$
  \end{itemize}
\item Perpetual put: $L^*=\frac{2rK}{2r+\sigma^2}=30$, \ $V(36)=5.79$
\item Finite $T$: Crank--Nicolson in $\log S$ + penalty for $V\ge g$
\item \textbf{Optional stopping:} for a martingale no rule beats holding,
  so memory can only help when increments are dependent.
\end{itemize}
\end{column}
\begin{column}{0.42\textwidth}
\small Boundary, $S_0=36$, $\sigma=0.2$, $T=1$:\\[2pt]
\begin{tabular}{rr}
\toprule
$t$ & $L(t)$\\
\midrule
0.00 & 32.94\\
0.25 & 33.34\\
0.50 & 33.99\\
0.75 & 35.06\\
0.95 & 37.05\\
$T^-$ & 40.00\\
\bottomrule
\end{tabular}
\end{column}
\end{columns}
\end{frame}

\begin{frame}{Ground truth: tree and PDE agree to $4\times10^{-4}$}
\centering\scriptsize
\begin{tabular}{rrr r rr rr}
\toprule
$S_0$ & $\sigma$ & $T$ & European & \multicolumn{2}{c}{Bermudan (50/yr)} & \multicolumn{2}{c}{American}\\
 & & & & tree & PDE & tree & PDE\\
\midrule
\csname @@input\endcsname ../results/benchmarks_table.tex
\bottomrule
\end{tabular}
\end{frame}

\begin{frame}{Baselines (P2)}
\begin{itemize}
\item \textbf{Longstaff--Schwartz:} regress realised cash flows on $[1,x,x^2]$ of today's
  features, in-the-money paths only, backwards in time.
\item \textbf{Deep Optimal Stopping:} one small network per date, trained backwards,
  each warm-started from the next date.
\item \textbf{Fitted Q-iteration:} one network (MLP or GRU) for every date, regressed on
  bootstrapped targets $\max(g_{j+1},K_{j+1})$.
\item All see the same features; all prices measured out of sample (lower bounds).
\end{itemize}
\end{frame}

% =============================================================== P1: OSPG
\section{OSPG}

\begin{frame}{OSPG: learn the stopping rule directly}
\begin{itemize}
\item Policy outputs $\varphi_j=\mathbb P(\text{stop at }j\mid\text{not yet},\ \mathcal F_j)$
  from the history (GRU) or today only (MLP).
\item Stopping does not change the path, so the stopping-time law is exact:
  \[\psi_j=\varphi_j\prod_{n<j}(1-\varphi_n),\qquad J(\theta)=\mathbb E\Big[\sum_j\psi_jg_j\Big].\]
\item Policy gradient with a constant baseline:
  $\nabla_\theta J=\mathbb E\big[\sum_j(g_j-b)\psi_j\nabla_\theta\log\psi_j\big]$.
\item No rollouts, \textbf{no backward recursion}: no bias carried from date to date.
\item Implicit Snell comparison: $\varphi_j$ rises when
  $\dfrac{\psi_jg_j}{\sum_{n>j}\psi_ng_n}>\dfrac{\varphi_j}{1-\varphi_j}$.
\end{itemize}
\end{frame}

\begin{frame}{What it took to make it work}
\begin{itemize}
\item \textbf{Start mostly waiting} (output bias $-2$): from $\varphi\approx\frac12$ the
  put policy collapsed to ``stop now'' (4.0 instead of 4.478).
\item \textbf{Small penalty on the logits}: the gradient is proportional to $\varphi$;
  without it the GRU stalled at ``never stop'' on the 50-asset max-call.
\item \textbf{Date 0 decided separately}: every simulated path starts in the same state.
\item Batch norm in the MLP (as in the paper), using training statistics at decision time.
\item Tests check the loss gradient equals the policy gradient, and that no method
  looks ahead.
\end{itemize}
\end{frame}

\begin{frame}{Replication: OSPG matches the exact prices}
\centering\scriptsize
% generated by experiments/summarize.py: method & mean (sd) per problem
\begin{tabular}{lrrrrr}
\toprule
Method & Put & Max-call $d{=}20$ & Max-call $d{=}50$ & Geo-call $d{=}7$ & Geo-call $d{=}20$ \\
\midrule
LSM & 4.458 (0.004) & 51.69 (0.20) & 69.83 (0.32) & 10.169 (0.012) & 9.899 (0.011) \\
DOS & 4.465 (0.007) & 51.69 (0.16) & 69.72 (0.29) & -- & -- \\
FQI-MLP & 4.467 (0.011) & 51.68 (0.12) & 69.67 (0.21) & 10.145 (0.068) & 9.874 (0.056) \\
FQI-GRU & -- & 51.73 (0.14) & 69.81 (0.37) & -- & -- \\
OSPG-MLP & 4.473 (0.008) & 51.80 (0.20) & 69.93 (0.28) & 10.255 (0.046) & 9.989 (0.017) \\
OSPG-GRU & 4.465 (0.015) & 51.72 (0.21) & 69.87 (0.25) & 10.237 (0.065) & 9.941 (0.030) \\
\midrule
Reference & 4.478 & 51.67 & 69.74 & 10.260 & 10.033 \\
Paper, DNN-OSPG & -- & 51.02 & 69.06 & 10.252 & 10.010 \\
\bottomrule
\end{tabular}


\vspace{6pt}\normalsize
\begin{itemize}\small
\item All methods within 1.6\% of the exact price; \textbf{OSPG-MLP closest}:
  $-0.12\%$ (put), $-0.05\%$ and $-0.44\%$ (geometric call), vs $-0.9$ to $-1.6\%$ for LS and FQI.
\item Paired on the same paths: OSPG-GRU beats LS and FQI on the geometric call ($z\ge2.6$),
  but trails OSPG-MLP by 0.1--0.5\%: memory costs a little when there is none.
\item The paper's max-call values are $\approx0.8$ below Becker et al.'s lower bound; we match Becker.
\end{itemize}

\end{frame}
   % (P2) optimal stopping, free boundary, baselines; (P1) OSPG, replication
% Slide block for P3 (fBm and rough volatility). Frames only; included by slides.tex.
\begin{frame}{Rough volatility (P3)}
\begin{itemize}
\item To do: fBm and rough Bergomi, fBm stopping table, American put under rough vol.
\end{itemize}
\end{frame}
      % (P3) fBm and rough volatility
% Slide block for P4 (data and backtest harness). Frames only; included by slides.tex.
\begin{frame}{Real market data and reproducible windows (P4)}
\begin{itemize}
\item Adjusted equity closes: 50 Indian and 100 usable US tickers; Binance spot: eight crypto assets.
\item Frozen 2015--2024 equity history; crypto requested from 2017. Retrieval-date universes have selection bias.
\item Audited cleaning: no future filling; missing windows excluded; large moves flagged and retained.
\item 30-calendar-day deadline; five-observation entry stride; 72,048 windows before fold selection.
\item Causal return, drawdown, volatility, lagged returns, remaining time; lagged local index/VIX context.
\end{itemize}
\end{frame}

\begin{frame}{Execution, costs and fair comparisons (P4)}
\[
r_j=\frac{S_j}{S_0}\frac{1-c_s}{1+c_b}-1-fd_j
\]
\begin{itemize}
\item Daily-close reference-price experiment; guaranteed same-close execution is not assumed.
\item Fees, half-spread and slippage on both legs; configurable research scenarios and cost sensitivity.
\item Hold to deadline; fixed $+5\%/-3\%$ profit/loss rule; $5\%$ trailing stop.
\item Expanding 36-month train / 6-month validation / 6-month test folds; purge overlapping outcomes.
\item Linear and exponential-utility rewards; actual net returns retained for both.
\end{itemize}
\end{frame}

\begin{frame}{Validation and the statistics handoff (P4)}
\begin{itemize}
\item 71 tests passed, including 12 P4 accounting, causality, split and integration tests.
\item Real-data smoke verified LSM, DOS, FQI-MLP/GRU and OSPG-MLP/GRU interfaces.
\item Saved asset/window IDs, entry/deadline/decision dates, stopping indices and actual exit dates.
\item Full rule runs: 12 walk-forward folds per equity market; seven crypto folds.
\item P5 performs final training and date-block inference: overlapping trades are dependent and are not portfolio returns.
\end{itemize}
\end{frame}
      % (P4) data and backtest harness
% Slide block for P5 (backtest results). Frames only; included by slides.tex.
\begin{frame}{Backtest results (P5)}
\begin{itemize}
\item To do: results by asset class, bootstrap intervals, regimes, memory vs Hurst.
\end{itemize}
\end{frame}
      % (P5) backtest results

\end{document}
